% GENERATED FILE. Edit canonical lessons, then run npm run generate:books.
% canonical-source-hash: fnv1a64:555aec754035d9b6
% canonical-lessons: LA-C539-alea, LA-C539-fritillus, LA-C539-latrunculus, LA-C539-lusus, LA-C539-luctor

\chapter{Dice and board games}
\label{ch:la-dice-and-board-games}
\hlchaptermodality{539}

\begin{chapteropening}
\textbf{By the end of this chapter:} \emph{I can talk about dice and board games.}

Complete the last lesson of chapter 539: I can talk about dice and board games.
\end{chapteropening}

\section[ālea, āleae]{ālea, āleae — dice-playing, a game of chance}

\label{lesson:LA-C539-alea}



\begin{quote}
Before the new one: say the Latin for to relax, then the Latin for relaxation.
\end{quote}

\subsection*{You'll want to know: ālea, āleae}
\textbf{ālea, āleae} — \textquotedblleft{}dice-playing, a game of chance\textquotedblright{}.

\begin{quote}
Mīlitēs in castrīs āleā lūdunt.
\end{quote}

\begin{quote}
\emph{The soldiers play dice in the camp.}
\end{quote}

Feminine, first declension. Gambling at dice was forbidden except at festival time. \emph{Ālea iacta est} — the die is cast.

\begin{cousinweb}
\textbf{ālea} is a word of uncertain origin; Lewis \& Short reports a guess linking it with a word for bone, but nothing is settled. English \emph{aleatory}, depending on chance, comes from it.
\end{cousinweb}

\begin{grammarlens}[title={What you've built}]
One more word for telling, explaining and describing at B1.
\end{grammarlens}

\subsection*{Your turn}
\begin{itemize}
\raggedright
  \item \emph{Say it:} \emph{ālea, āleae}
  \item \emph{Say it:} \emph{ālea, āleae}, once more
  \item \emph{Say it:} a sentence of your own with \emph{ālea, āleae}
  \item \emph{Recall:} say the Latin for to relax, then the Latin for relaxation, then say \emph{ālea, āleae} again
\end{itemize}

\begin{tcolorbox}[breakable,colback=teal!4,colframe=teal!35!black,title={Before you move on}]
What does ālea, āleae mean? (\textquotedblleft{}Dice-playing, a game of chance\textquotedblright{}.) Say it once more.
\end{tcolorbox}
\section[fritillus, fritillī]{fritillus, fritillī — a dice-box}

\label{lesson:LA-C539-fritillus}



\begin{quote}
Before the new one: say the Latin for relaxation, then the Latin for dice-playing.
\end{quote}

\subsection*{You'll want to know: fritillus, fritillī}
\textbf{fritillus, fritillī} — \textquotedblleft{}a dice-box\textquotedblright{}.

\begin{quote}
Senex tālōs ē fritillō iacit.
\end{quote}

\begin{quote}
\emph{The old man throws the knucklebones from the dice-box.}
\end{quote}

Masculine, second declension. A little cup for shaking and throwing the dice.

\begin{cousinweb}
\textbf{fritillus} is a native Latin noun whose origin is not settled (Lewis \& Short gives none). The flower \emph{fritillary} is named after it.
\end{cousinweb}

\begin{grammarlens}[title={What you've built}]
One more word for telling, explaining and describing at B1.
\end{grammarlens}

\subsection*{Your turn}
\begin{itemize}
\raggedright
  \item \emph{Say it:} \emph{fritillus, fritillī}
  \item \emph{Say it:} \emph{fritillus, fritillī}, once more
  \item \emph{Say it:} a sentence of your own with \emph{fritillus, fritillī}
  \item \emph{Recall:} say the Latin for relaxation, then the Latin for dice-playing, then say \emph{fritillus, fritillī} again
\end{itemize}

\begin{tcolorbox}[breakable,colback=teal!4,colframe=teal!35!black,title={Before you move on}]
What does fritillus, fritillī mean? (\textquotedblleft{}A dice-box\textquotedblright{}.) Say it once more.
\end{tcolorbox}
\section[latrunculus, latrunculī]{latrunculus, latrunculī — a piece in a board game; a robber}

\label{lesson:LA-C539-latrunculus}



\begin{quote}
Before the new one: say the Latin for dice-playing, then the Latin for a dice-box.
\end{quote}

\subsection*{You'll want to know: latrunculus, latrunculī}
\textbf{latrunculus, latrunculī} — \textquotedblleft{}a piece in a board game; a robber\textquotedblright{}.

\begin{quote}
Duo amīcī latrunculīs lūdunt.
\end{quote}

\begin{quote}
\emph{Two friends play the board game of soldiers.}
\end{quote}

Masculine, second declension. First a little robber or hired soldier, then a piece in a war game played on a board.

\begin{cousinweb}
\textbf{latrunculus} is a diminutive of \emph{latrō} \textquotedblleft{}robber, hired soldier\textquotedblright{} (Lewis \& Short): the pieces of the game were little soldiers.
\end{cousinweb}

\begin{grammarlens}[title={What you've built}]
One more word for telling, explaining and describing at B1.
\end{grammarlens}

\subsection*{Your turn}
\begin{itemize}
\raggedright
  \item \emph{Say it:} \emph{latrunculus, latrunculī}
  \item \emph{Say it:} \emph{latrunculus, latrunculī}, once more
  \item \emph{Say it:} a sentence of your own with \emph{latrunculus, latrunculī}
  \item \emph{Recall:} say the Latin for dice-playing, then the Latin for a dice-box, then say \emph{latrunculus, latrunculī} again
\end{itemize}

\begin{tcolorbox}[breakable,colback=teal!4,colframe=teal!35!black,title={Before you move on}]
What does latrunculus, latrunculī mean? (\textquotedblleft{}A piece in a board game; a robber\textquotedblright{}.) Say it once more.
\end{tcolorbox}
\section[lūsus, lūsūs]{lūsus, lūsūs — play, a game}

\label{lesson:LA-C539-lusus}



\begin{quote}
Before the new one: say the Latin for a dice-box, then the Latin for a piece in a board game.
\end{quote}

\subsection*{You'll want to know: lūsus, lūsūs}
\textbf{lūsus, lūsūs} — \textquotedblleft{}play, a game\textquotedblright{}.

\begin{quote}
Lūsus puerōrum in viā clāmōsus est.
\end{quote}

\begin{quote}
\emph{The children's play in the street is noisy.}
\end{quote}

Masculine, fourth declension, from \emph{lūdō} — to play.

\begin{cousinweb}
\textbf{lūsus} is the action noun of \emph{lūdō} \textquotedblleft{}to play\textquotedblright{} (Lewis \& Short; EtymDB). English \emph{illusion} comes from the family, through \emph{illūdō} \textquotedblleft{}to mock\textquotedblright{}.
\end{cousinweb}

\begin{grammarlens}[title={What you've built}]
One more word for telling, explaining and describing at B1.
\end{grammarlens}

\subsection*{Your turn}
\begin{itemize}
\raggedright
  \item \emph{Say it:} \emph{lūsus, lūsūs}
  \item \emph{Say it:} \emph{lūsus, lūsūs}, once more
  \item \emph{Say it:} a sentence of your own with \emph{lūsus, lūsūs}
  \item \emph{Recall:} say the Latin for a dice-box, then the Latin for a piece in a board game, then say \emph{lūsus, lūsūs} again
\end{itemize}

\begin{tcolorbox}[breakable,colback=teal!4,colframe=teal!35!black,title={Before you move on}]
What does lūsus, lūsūs mean? (\textquotedblleft{}Play, a game\textquotedblright{}.) Say it once more.
\end{tcolorbox}
\section[luctor, luctārī]{luctor, luctārī — to wrestle, struggle}

\label{lesson:LA-C539-luctor}



\begin{quote}
Before the new one: say the Latin for a piece in a board game, then the Latin for play.
\end{quote}

\subsection*{You'll want to know: luctor, luctārī}
\textbf{luctor, luctārī} — \textquotedblleft{}to wrestle, struggle\textquotedblright{}.

\begin{quote}
Iuvenēs in campō luctantur.
\end{quote}

\begin{quote}
\emph{The young men wrestle on the field.}
\end{quote}

Deponent: \emph{luctor, luctārī, luctātus sum}. Also to struggle against something. English \emph{reluctant} comes from \emph{reluctor}.

\begin{cousinweb}
\textbf{luctor} is built on \emph{lucta} \textquotedblleft{}wrestling\textquotedblright{} (Lewis \& Short); EtymDB links it with a Proto-Indo-European root meaning \textquotedblleft{}to bend, twist\textquotedblright{}. English \emph{reluctant} comes from the compound \emph{reluctor}.
\end{cousinweb}

\begin{grammarlens}[title={What you've built}]
One more word for telling, explaining and describing at B1.
\end{grammarlens}

\subsection*{Your turn}
\begin{itemize}
\raggedright
  \item \emph{Say it:} \emph{luctor, luctārī}
  \item \emph{Say it:} \emph{luctor, luctārī}, once more
  \item \emph{Say it:} a sentence of your own with \emph{luctor, luctārī}
  \item \emph{Recall:} say the Latin for a piece in a board game, then the Latin for play, then say \emph{luctor, luctārī} again
\end{itemize}

\begin{tcolorbox}[breakable,colback=teal!4,colframe=teal!35!black,title={Before you move on}]
What does luctor, luctārī mean? (\textquotedblleft{}To wrestle, struggle\textquotedblright{}.) Say it once more.
\end{tcolorbox}
